\section{Introduction}
\label{sec:intro}

\sysname is an infrastructure-observability platform. Its agents run system and
network monitoring on the hosts they watch, reporting CPU, memory, interface
counters, link health, and service availability on a fixed cadence. Across a fleet
the number of distinct metric series runs into the millions. Operators want to know
quickly when any of those series misbehaves. The obvious way to find out, shipping
every sample to a central service and analyzing it there, makes the central tier
pay the full cost of detection and grow that cost with the fleet. We take the
opposite approach: we detect on the agent, as the data is produced, with a detector
framed in the language of computational causality.

\paragraph{Why a causal engine.}
We model each series as a \emph{causaloid} from \dc, a small unit of causal
inference evaluated against an explicit context, here the series' own recent window.
A rolling $z$-score with confirmation is not itself new. What the causal formulation
buys us is the shape of the engine around it. The unit is small and tested on its
own. Reasoning is stateless over a bounded context, which makes a series cheap to
checkpoint and restart. The same reasoner produces identical verdicts whether it
runs on the agent or in the core, so we write it once and trust it everywhere.
Because the detector is already expressed in \dc, it also has a direct path to the
framework's intervention step, which would let the same engine take guarded
corrective action on the host instead of only reporting. We do not take that step
here, but nothing in the design forecloses it.

\paragraph{Thesis.}
Detection can live on every agent if the reasoner is cheap enough to run on each
sample as it arrives. \dc{} gives us that reasoner. The result is an engine that
flags anomalous metrics across the fleet in real time, inside a small CPU and
memory budget, without centralizing raw telemetry.

\paragraph{Contributions.}
\begin{itemize}
  \item A deployment model for anomaly detection: a real-time, per-series detector
        cheap and self-contained enough to run on every monitoring agent, so the
        central tier no longer pays for detection or grows that cost with the fleet.
  \item A causaloid-based formulation of that detector on \dc{} (a
        \texttt{CausalFlow} pipeline over a bounded-window context, with
        baseline-integrity admission and confirmation hysteresis), and an account of
        what the formulation buys: composability, stateless checkpointable state,
        on-agent and in-core verdict parity, and an on-ramp to causal interventions.
  \item A shared Rust core that computes edge verdicts identically to the reference
        implementation by construction, checked by a parity test, with correct
        handling of gauges and rate-normalized counters.
  \item Microbenchmarks on commodity hardware that separate the reasoner's
        arithmetic from the per-series engine, showing detection keeps up in real
        time well inside a fixed resource budget, with an honest account of what we
        have not yet measured.
\end{itemize}
